\begin{table}[!h]
\centering
\caption{\label{tab:descriptive}Descriptive Statistics in 1990}
\centering
\resizebox{\ifdim\width>\linewidth\linewidth\else\width\fi}{!}{
\begin{tabular}[t]{>{\raggedright\arraybackslash}p{7cm}ccc}
\toprule
\multicolumn{1}{c}{ } & \multicolumn{3}{c}{Groups} \\
\cmidrule(l{3pt}r{3pt}){2-4}
Variable & Never Treated & Treated after 1995 & Treated in 1995\\
\midrule
\addlinespace[0.3em]
\multicolumn{4}{l}{\textbf{Past elections}}\\
\hspace{1em}Vote share for FN in 1988 & \makecell{14.21 \\ (6.15)} & \makecell{11.95 \\ (5.55)} & \makecell{10.67 \\ (6.02)}\\
\hspace{1em}FN vote share in 1995 & \makecell{16.60 \\ (6.41)} & \makecell{14.03 \\ (6.21)} & \makecell{11.64 \\ (6.63)}\\
\addlinespace[0.3em]
\multicolumn{4}{l}{\textbf{Employment}}\\
\hspace{1em}Unemployed (\%) & \makecell{8.44 \\ (5.61)} & \makecell{9.01 \\ (7.19)} & \makecell{8.78 \\ (9.18)}\\
\hspace{1em}In the labor force (\%) & \makecell{23.59 \\ (30.40)} & \makecell{24.14 \\ (17.96)} & \makecell{25.15 \\ (17.69)}\\
\hspace{1em}Agriculture (\%) & \makecell{8.25 \\ (11.25)} & \makecell{17.96 \\ (17.57)} & \makecell{25.76 \\ (22.57)}\\
\hspace{1em}Independant (\%) & \makecell{8.33 \\ (6.26)} & \makecell{8.37 \\ (8.24)} & \makecell{9.04 \\ (10.53)}\\
\hspace{1em}Intermediate occupations (\%) & \makecell{19.24 \\ (8.81)} & \makecell{14.63 \\ (10.34)} & \makecell{13.53 \\ (12.15)}\\
\hspace{1em}Clerical (\%) & \makecell{22.47 \\ (8.52)} & \makecell{18.90 \\ (10.82)} & \makecell{17.96 \\ (13.15)}\\
\hspace{1em}Manual (\%) & \makecell{33.65 \\ (13.57)} & \makecell{35.33 \\ (15.84)} & \makecell{29.26 \\ (17.77)}\\
\addlinespace[0.3em]
\multicolumn{4}{l}{\textbf{Demographics}}\\
\hspace{1em}Population & \makecell{2795.31 \\ (20421.44)} & \makecell{699.03 \\ (1579.92)} & \makecell{403.39 \\ (846.54)}\\
\hspace{1em}Foreigners (\%) & \makecell{2.60 \\ (3.52)} & \makecell{1.60 \\ (2.64)} & \makecell{1.86 \\ (2.92)}\\
\hspace{1em}Ages 20-40 (\%), men & \makecell{18.68 \\ (3.94)} & \makecell{17.86 \\ (5.20)} & \makecell{16.78 \\ (6.39)}\\
\hspace{1em}Ages 20-40 (\%), women & \makecell{17.67 \\ (3.71)} & \makecell{16.03 \\ (4.68)} & \makecell{14.48 \\ (5.58)}\\
\hspace{1em}Age ratio young/old (\%) & \makecell{80.41 \\ (22.41)} & \makecell{69.21 \\ (22.31)} & \makecell{58.28 \\ (24.24)}\\
\hspace{1em}Population density & \makecell{2.57 \\ (9.00)} & \makecell{0.51 \\ (1.20)} & \makecell{0.25 \\ (0.47)}\\
\hspace{1em}Population change in p.p. 1980-1990 & \makecell{0.21 \\ (0.34)} & \makecell{0.05 \\ (0.19)} & \makecell{-0.02 \\ (0.21)}\\
\hspace{1em}Vacant housing (\%) & \makecell{6.62 \\ (3.90)} & \makecell{8.87 \\ (4.79)} & \makecell{10.21 \\ (5.95)}\\
\hspace{1em}OPI per 1,000 inhabitants & \makecell{2.78 \\ (0.80)} & \makecell{3.02 \\ (0.73)} & \makecell{3.26 \\ (0.79)}\\
\hspace{1em}Taxable income per capita  (log) & \makecell{10.57 \\ (0.25)} & \makecell{10.43 \\ (0.23)} & \makecell{10.35 \\ (0.23)}\\
\bottomrule
\end{tabular}}
\end{table}
\clearpage
\begin{table}[!h]
\centering
\caption{Descriptive Statistics in 1990 (continued)}
\centering
\resizebox{\ifdim\width>\linewidth\linewidth\else\width\fi}{!}{
\begin{threeparttable}
\begin{tabular}[t]{>{\raggedright\arraybackslash}p{7cm}ccc}
\toprule
\multicolumn{1}{c}{ } & \multicolumn{3}{c}{Groups} \\
\cmidrule(l{3pt}r{3pt}){2-4}
Variable & Never Treated & Treated after 1995 & Treated in 1995\\
\midrule
\addlinespace[0.3em]
\multicolumn{4}{l}{\textbf{Education}}\\
\hspace{1em}No diploma (\%) & \makecell{21.30 \\ (8.76)} & \makecell{25.93 \\ (10.08)} & \makecell{26.46 \\ (11.76)}\\
\hspace{1em}Academic (\%) & \makecell{5.61 \\ (4.08)} & \makecell{3.96 \\ (3.67)} & \makecell{4.13 \\ (4.48)}\\
\hspace{1em}Highschool (\%) & \makecell{6.82 \\ (3.43)} & \makecell{5.82 \\ (4.01)} & \makecell{6.22 \\ (5.06)}\\
\hspace{1em}Technical (\%) & \makecell{15.98 \\ (4.98)} & \makecell{14.10 \\ (5.69)} & \makecell{13.66 \\ (7.02)}\\
\addlinespace[0.3em]
\multicolumn{4}{l}{\textbf{Geography}}\\
\hspace{1em}Altitude & \makecell{4.91 \\ (0.99)} & \makecell{5.03 \\ (0.84)} & \makecell{5.76 \\ (0.76)}\\
\hspace{1em}Distance to closest agglomeration in meters (log) & \makecell{9.48 \\ (1.51)} & \makecell{10.20 \\ (0.52)} & \makecell{10.44 \\ (0.42)}\\
\hspace{1em}Area in km2 (log) & \makecell{6.87 \\ (0.80)} & \makecell{7.02 \\ (0.75)} & \makecell{7.22 \\ (0.76)}\\
\hspace{1em}Observations & 17047 & 6618 & 10898\\
\bottomrule
\end{tabular}
\begin{tablenotes}
\item \textit{Note: } 
\item Cells report the mean (first line) and standard deviation (in parentheses, second line). Some variables are taken from nearby years: taxable income (1994) and young/old ratio (1995).
\item[a] In the socio-economic database assembled by Julia Cag\'e and Thomas Piketty (2023), the average income per municipality is defined as the total income reported on tax declarations (before any deductions or allowances) divided by the total number of inhabitants (including children). Source: \url{https://www.unehistoireduconflitpolitique.fr/glossaire.html}, the website associated with the book by Julia Cag\'e and Thomas Piketty (2023): \textit{Une histoire du conflit politique. \'{E}lections et in\'{e}galit\'{e}s sociales en France, 1789--2022}, Paris, Le Seuil.
\end{tablenotes}
\end{threeparttable}}
\end{table}
